\chapter{Teilschaltungen}
\label{kap:format-teilschaltungen}

\section{Definition und Instanz}

Wiederverwendbare Teilschaltungen stehen unter \code{subcircuits} und
werden mit dem Bauteil \code{X} eingesetzt:

\begin{json}
"subcircuits": {
  "divider": {
    "ports": ["in", "out"],
    "parameters": { "ratio": 0.5, "rtotal": "10k" },
    "components": [
      { "type": "R", "name": "R1", "nodes": ["in", "out"], "value": "(1 - ratio)*rtotal" },
      { "type": "R", "name": "R2", "nodes": ["out", "0"],  "value": "ratio*rtotal" }
    ]
  }
},
"components": [
  { "type": "X", "name": "X1", "subcircuit": "divider", "nodes": ["a", "b"] },
  { "type": "X", "name": "X2", "subcircuit": "divider", "nodes": ["a", "c"],
    "parameters": { "ratio": 0.25 } }
]
\end{json}

\begin{description}[style=nextline,leftmargin=2em]
  \item[Knoten]
    Die Namen in \code{ports} stehen für die außen angeschlossenen Knoten,
    in der Reihenfolge der \code{nodes} der Instanz. Alle anderen Knoten
    sind lokal. Masse (\code{"0"}) ist überall dieselbe.
  \item[Parameter]
    \code{parameters} der Definition enthält Standardwerte; sie dürfen sich
    auf globale Parameter und aufeinander beziehen. Die Instanz
    überschreibt einzelne Werte mit \code{"parameters": {...}}.
  \item[Inhalt]
    Alle Bauteile, auch weitere Instanzen (Schachtelung bis 20 Ebenen),
    dazu eigene \code{models}, \code{variables}, \code{equations} und
    \code{events}. Ein Bauteil mit thermischem Zustand oder eigener
    Kennlinie lässt sich so als Modell kapseln und mehrfach verwenden.
  \item[Plan]
    Der Editor legt den Plan der Teilschaltung unter \code{"schematic"} in
    ihrer Definition ab.
\end{description}

\section{Namen von außen}

Was in der Instanz \code{X1} liegt, heißt außen \code{X1.<Name>}, bei
Schachtelung \code{X1.U.<Name>}:

\begin{itemize}
  \item Knotenspannungen: \code{v(X1.mid)},
  \item Ströme: \code{i(X1.V1)},
  \item Variablen: \code{X1.T},
  \item Quellen eines DC-Sweeps: \code{"source": "X1.V1"}.
\end{itemize}

Innerhalb einer Teilschaltung beziehen sich \code{i(...)} und
\code{control} auf Bauteile derselben Teilschaltung. In der Standardausgabe
(ohne \code{outputs}) erscheint das Innere von Instanzen nicht; innere
Größen müssen ausdrücklich unter \code{outputs} stehen.

\section{Teilschaltungen in SPICE}
\label{sec:spice-subckt}

Modelle der Hersteller lassen sich unverändert einsetzen: Statt eines
Objekts steht der SPICE-Text, als eine Zeichenkette oder als Liste von
Zeilen.

\begin{json}
"subcircuits": {
  "OPX": [
    ".SUBCKT OPX inp inn out vp vn",
    "Rin inp inn 1meg",
    "Egain mid 0 POLY(1) inp inn 0 1e5",
    "Rout mid out 10",
    "Dhi out vp DCLAMP",
    "Dlo vn out DCLAMP",
    ".MODEL DCLAMP D(IS=1e-14)",
    ".ENDS"
  ]
},
"components": [
  { "type": "X", "name": "U1", "subcircuit": "OPX", "nodes": ["in", "fb", "out", "vp", "vn"] }
]
\end{json}

Die Modellbibliothek des Editors erzeugt genau diese Form
(Kapitel~\ref{kap:editor-bibliothek}).

\subsection*{Was gelesen wird}

\begin{itemize}
  \item Bauteile \code{R}, \code{C}, \code{L}, \code{K}, \code{V} und
    \code{I} (DC, AC, SIN, PULSE, PWL), \code{E} und \code{G} (linear,
    \code{POLY(n)}, \code{VALUE = {...}}, \code{TABLE {...} = (x, y) ...}),
    \code{F} und \code{H} (linear, \code{POLY(n)}), \code{B} (\code{V=}
    oder \code{I=}), \code{D}, \code{Q}, \code{M}, \code{J}, \code{S},
    \code{T} und \code{X}.
  \item \code{.param}, \code{PARAMS:} an \code{.subckt} und an Instanzen,
    \code{.model} – auch mit Ausdrücken, ausgewertet mit den Parametern der
    Teilschaltung.
  \item Formen von PSpice: Widerstände mit Modell (\code{R1 a b RMOD 1k};
    \code{R=} im Modell skaliert, \code{T_ABS=-273.15} macht den Widerstand
    rauschfrei); Schalter mit \code{VSWITCH(RON ROFF VON VOFF)} mit dem
    stetigen Übergang von PSpice zwischen \code{VOFF} und \code{VON}
    (intern eine Verhaltensquelle); Schalter mit \code{SW(VT VH)} werden
    esims Schalter mit Hysterese.
  \item Ausdrücke in \code{{...}}: Zahlen mit SPICE-Endungen (\code{M} ist
    Milli, \code{meg} Mega), \code{**}, \code{a ? b : c}, \code{limit},
    \code{pwr}, \code{pwrs}, \code{uramp}, \code{u}, \code{stp},
    \code{sgn}.
\end{itemize}

Modelle und weitere \code{.subckt} im Text gehören zu dieser
Teilschaltung und heißen intern \code{OPX/NAME}; sie stoßen nicht mit
Modellen der Datei zusammen. Was der Text nicht definiert, sucht esim in
der Datei.

\subsection*{Was nicht geht}

Nicht Unterstütztes (z.\,B. \code{LAPLACE}, \code{.include}, unbekannte
Bauteile) ist ein Fehler, der die Zeile im Text nennt. Ohne Wirkung
bleiben, mit Warnung, etwa die Fläche von Dioden und der Substratknoten
von Transistoren.

\subsection*{Erprobt}

\begin{itemize}
  \item Das Makromodell eines verbreiteten Doppel-Operationsverstärkers
    (PSpice, 18 Teilschaltungen, 227 Bauteile): DC-Kennlinie und
    Sprungantwort stimmen mit ngspice im PSpice-Modus überein (DC auf 7
    Stellen, Zeitverlauf auf 0,3\,mV).
  \item Ein kleiner P-Kanal-MOSFET als Level-3-Modell eines Herstellers:
    Kennlinien wie ngspice auf 6 Stellen, Schaltflanke auf 17\,mV bei 5\,V
    Hub.
  \item \datei{tests/spice_subckt_ngspice.py} vergleicht eine SPICE-
    Teilschaltung mit allen Formen dieses Abschnitts gegen ngspice.
\end{itemize}
